\addcontentsline{toc}{subsubsection}{Solution 8. The Multiplicative Mirror}

\begin{dlsummary}[title={\bfseries Solution 8. The Multiplicative Mirror},breakable]

\begin{proof}
The reciprocal condition gives $f(1)=1$ and $f(1/x)=1/f(x)$. The
inequalities become
\[
\frac1{f(x+y)}\geq\frac1{f(x)}+\frac1{f(y)},\qquad
f(xy)f(x/y)\geq f(x)^2. \tag{1}
\]
Replacing $x$ by $1/x$ in the second inequality and using reciprocity
gives the reverse inequality. Hence
\[
f(xy)f(x/y)=f(x)^2.
\]
Applying this to $(u,v)$ and $(v,u)$ and multiplying yields
$f(uv)^2=f(u)^2f(v)^2$, so positivity gives $f(uv)=f(u)f(v)$.

Set $g=1/f$. Then $g$ is positive, increasing, superadditive,
multiplicative, and $g(1)=1$. The reciprocal condition also gives
\[
g(1/x)=\frac1{f(1/x)}=f(x)=\frac1{g(x)}.
\]
In particular $g(2)\geq2$. Moreover, using this identity with
$x=f(2)$,
\[
g(g(2))=\frac1{f(1/f(2))}=f(f(2))\leq2.
\]
Since $g(2)\geq2$ and $g$ is increasing,
$2\geq g(g(2))\geq g(2)\geq2$, so $g(2)=2$.

Multiplicativity gives $g(2^q)=2^q$ for every rational $q$. These
arguments form a dense subset of $(0,\infty)$; monotonicity therefore
forces $g(x)=x$ everywhere. Thus $f(x)=1/x$. Direct substitution
shows this function satisfies all conditions with equality.
\end{proof}

\end{dlsummary}